\scanpage{16}
A - 21 -\hfill n\textsuperscript{o} 413

de V dans N comme étant un élément de $\operatorname{Hom}(U(\underline V);N)$ et noter $\widehat P(\underline V;N)$ le module ainsi défini, $\widehat P_+(\underline V;N)$ désignant le sous-module des applications linéaires nulles sur $U_0(\underline V)$. Ceci ne peut être gênant puisque dans le cas d'un module libre, il y a un isomorphisme canonique de l'un sur l'autre des deux modules notés $\widehat P_+(\underline V;N)$.
Si $N=K$, on identifiera $\widehat P(\underline V;K)$ et $\widehat P(\underline V)$.

Soient $\underline V$ et $\underline W$ deux noyaux formels. Nous allons caractériser les applications linéaires de $U(\underline V)$ dans $U(\underline W)$ qui proviennent d'applications formelles de $\underline V$ dans $\underline W$:

PROPOSITION 8.--- \uline{Pour qu'une application linéaire} $\psi$ \uline{de} $U(\underline V)$ \uline{dans} $U(\underline W)$ \uline{soit de la forme} $\varphi_*$ \uline{avec} $\varphi\in\widehat P_+(\underline V;\underline W)$, \uline{il faut et il suffit que} $\psi(U_0(\underline V))=U_0(\underline W)$ \uline{et que} $\psi$ \uline{soit compatible avec les applications diagonales} (\uline{c'est-à-dire que}
\[
(\psi\otimes\psi)\circ\unsure{\Gamma|T}=\unsure{\Gamma|T}\circ\psi).
\]
Que les conditions soient nécessaires résulte de la prop. 6. La prop. 7 montre d'ailleurs que si $\varphi$ existe, elle est unique. Pour démontrer l'existence de $\varphi$, on se ramène en prenant des cartes formelles, au cas de noyaux formels définis par des modules libres M et N. Soit alors $(e_i)_{i\in I}$ une base de M, $(e'_j)_{j\in J}$ une base de M', et posons:
\[
(30)\qquad\psi(e_\alpha)=\sum a_{\alpha,\beta}\,e'_\beta\qquad(a_{\alpha,\beta}\in K).
\]
La condition $\psi(U_0(V))=U_0(W)$ entraîne
\[
(31)\qquad a_{0,\beta}=0\quad\text{pour }|\beta|\geq1,\quad a_{0,0}\ne0.
\]
La compatibilité de $\psi$ avec les applications diagonales donne:
\[
(32)\qquad a_{\alpha,\beta_1+\beta_2}=\sum_{\alpha_1+\alpha_2=\alpha}a_{\alpha_1,\beta_1}\,a_{\alpha_2,\beta_2}.
\]
\figure{\begin{tikzpicture}[scale=.6]\draw[line width=.7pt] (0,.5)--(.23,.17);\draw[line width=.7pt] (-.11,-.28)--(-.11,-.45);\fill (-.2,-.3) circle (.012);\end{tikzpicture}}
\mspage{15}
